\section{Methods}

\paragraph{Design and pre-registration.} The endpoints, cross-validation scheme,
baseline feature set, and the accepted null were frozen and git-tagged before any
label was inspected. The null hypothesis that SFI merely re-encodes fibrosis was
pre-registered as an accepted, reportable outcome. The section of the
pre-registration governing the clinical analysis was frozen before any feature was
placed beside the outcome column; that analysis was run once and is reported as it
came out.

\paragraph{Reporting guideline.} The clinical endpoint --- prediction of two-year
post-ablation arrhythmia recurrence in the 82 University of Washington (UW)
patients, the only analysis in this work scored against a real patient outcome ---
is reported in accordance with TRIPOD+AI, and a completed checklist accompanies
this article. Adherence is claimed for that endpoint alone: every other outcome
here is an electrophysiology simulator's reentry-inducibility verdict rather than a
patient outcome, so those analyses lie outside the statement's remit and no claim
of adherence is made for them, though they follow the same frozen protocol.

\paragraph{Cohorts.} Two public patient-derived left-atrial mesh deposits were
used. The Roney virtual cohort (Zenodo 5801337) contains 100 meshes and all 100
were analysed; fibrosis is a continuous fraction mapped from the shipped
image-intensity ratio. The UW/Boyle late-gadolinium-enhancement cohort (Dryad
10.5061/dryad.kkwh70sg0) contains 164 meshes, being 82 patients imaged before and
after ablation; the primary analysis used the 82 pre-ablation meshes only, because
post-ablation geometry encodes the delivered treatment. That release ships no
intensity ratios (withheld under its institutional review board), so its fibrosis
was derived from the per-element tissue tags and is two-valued. No subject was
excluded for any other reason, giving 182 analysed meshes with no attrition. Both
deposits are de-identified and carry anatomy only, with no age, sex, race,
ethnicity, or comorbidity variables. The per-patient two-year recurrence outcomes
were supplied on request by the UW cohort's authors, who confirmed they carry no
confidentiality restriction. Synthetic excitable networks (up to 100\,000 random
geometric graphs) were used only for the scale behaviour of the null and for
cross-medium transfer.

\paragraph{Conduction graph and index.} Nodes were mesh vertices and edges
triangle edges, with weight
$w_{ij}=[c_\perp+(c_\parallel-c_\perp)a_{ij}^{2}]\cdot\mathrm{clip}(1-\phi_{ij},\varepsilon,1)$,
where $a_{ij}$ is edge--fibre alignment, $\phi_{ij}$ mean endpoint fibrosis,
$c_\parallel=1.0$, $c_\perp=0.3$, and $\varepsilon=0.05$; $L=D-W$ has algebraic
connectivity $\lam_2$ with Fiedler vector $\phii_2$. Diffuse fibrosis-weighted
uncoupling $\Delta w$ (mean $\Delta w\approx0.36$, frozen in the pre-registration)
gives, to first order,
\begin{equation}
\SFI(R)\;\approx\;\sum_{(i,j)\in R}\mathbb{E}[\Delta w_{ij}]\,
(\phii_{2,i}-\phii_{2,j})^{2},
\qquad
\frac{\partial\lam_2}{\partial w_{ij}}=(\phii_{2,i}-\phii_{2,j})^{2},
\end{equation}
a classical identity. First-order, basis-independent subspace, and
exact-recompute variants were all computed on $\sim$2000-node coarsened graphs.

\paragraph{Comparator features and outcomes.} SFI was evaluated strictly as an
add-on to six published competitor features it had to beat: fibrosis burden,
fibrosis spatial entropy, fibrosis patch size, deterministic global min-cut
(Stoer--Wagner), percolation threshold, and $\lam_2$ alone. The in-silico outcome
was a monodomain Mitchell--Schaeffer inducibility verdict (self-sustained reentry
$\geq650$~ms after burst pacing at a cycle length of 150~ms from random sites),
42 of 182 meshes inducible; this is a simulator verdict and not clinical atrial
fibrillation. The clinical outcome was recurrence, defined as at least 30 seconds
of documented atrial fibrillation or flutter after a 90-day blanking period and
within two years; all 82 patients were followed the full two years, with no
censoring, and 48 had events.

\paragraph{Statistical analysis.} Two classifiers were used: a standardised,
penalised logistic regression whose regularisation was chosen in an inner
\code{GroupKFold} over $\{0.03,0.1,0.3,1.0,3.0\}$, and a deliberately untuned,
frozen gradient-boosted tree ensemble (120 trees, depth 2, learning rate 0.05,
subsample 0.8), the latter logged as a deviation from the pre-registered nested
selection. Cross-validation was \code{GroupKFold} with the group set to the
patient, so no patient appeared in both training and test folds, using
$k=\min(5,n_{\mathrm{groups}})$ folds; a naive shuffled-fold arm was reported
alongside so that leakage is visible. Out-of-fold predicted probabilities were
pooled and compared by the paired DeLong test in its fast midrank form; because
pooled scores come from $k$ fitted models rather than two fixed scoring functions,
these $p$ values are approximate. Intervals came from at least $10^{4}$ stratified
bootstrap resamples. The primary endpoint required \emph{both} grouped
$\dAUC\geq0.05$ and DeLong $p<0.05$. That effect-size gate was not reused for the
clinical arm, where power at $\dAUC=0.05$ with 82 patients and 48 events is 0.24;
the minimum detectable effect was instead fixed at $\dAUC\approx0.10$--$0.15$, and
the plan pre-committed that a baseline materially weaker than the source study's
published performance would make any increment confounded by baseline degradation.
Testing followed a fixed sequence halting at the first non-rejection. Full methods
appear in eMethods in Supplement 2.
